\uuid{LGAw}
\exo7id{2903}
\titre{exo7 2903}
\auteur{quercia}
\organisation{exo7}
\datecreate{2010-03-08}
\isIndication{false}
\isCorrection{true}
\chapitre{Dénombrement}
\sousChapitre{Binôme de Newton et combinaison}
\cpge{mpsi-mp2i.denombrement.a, pcsi.denombrement.a, ptsi.denombrement.a, mpsi-mp2i.denombrement.b, pcsi.denombrement.b, ptsi.denombrement.b, mpsi-mp2i.calcul-algebrique-trigonometrie.a, pcsi.calcul-algebrique-trigonometrie.a, ptsi.calcul-algebrique-trigonometrie.a, bcpst1.denombrement, bcpst1.methodes-calcul}
\module{Algèbre}
\niveau{L1}
\difficulte{}

\contenu{
\question{
Soit $E$ un ensemble fini de cardinal $n$.
Calculer $\sum_{A \subset E} \mathrm{Card}\,(A)$,
         $\sum_{A,B \subset E} \mathrm{Card}\,(A\cap B)$,
         $\sum_{A,B \subset E} \mathrm{Card}\,(A\cup B)$.
}
\reponse{
$n2^{n-1}$, $n4^{n-1}$, $3n4^{n-1}$.
}
}
